\chapter{Листинг точки входа приложения (cmd/analyzer/main.go)}

В приложении ниже приведён исходный код точки входа, отвечающей за инициализацию графического интерфейса и формирование основных вкладок комплекса.

\begin{lstlisting}[language=Go, caption={Исходный код файла cmd/analyzer/main.go}]
package main

import (
	"system-analyzer/pkg/gui"

	"fyne.io/fyne/v2"
	"fyne.io/fyne/v2/app"
	"fyne.io/fyne/v2/container"
)

// main is the entry point of the application.
func main() {
	// Create a new Fyne application.
	a := app.New()
	// Create the main application window with a title.
	w := a.NewWindow("Системный Анализатор")

	// Создаем контейнер с вкладками для различных панелей.
	tabs := container.NewAppTabs(
		container.NewTabItem("Мониторинг", container.NewScroll(gui.NewDashboardPanel())),
		container.NewTabItem("Процессор", container.NewScroll(gui.NewBenchmarkPanel("Процессор"))),
		container.NewTabItem("Память", container.NewScroll(gui.NewBenchmarkPanel("Память"))),
		container.NewTabItem("Математика", container.NewScroll(gui.NewBenchmarkPanel("Математика"))),
		container.NewTabItem("Криптография", container.NewScroll(gui.NewInteractiveCryptoPanel())),
	)

	// Set the tab container as the main window content.
	w.SetContent(tabs)
	// Set window to maximum available size
	w.Resize(fyne.NewSize(1600, 1000))
	w.CenterOnScreen()
	w.ShowAndRun()
}
\end{lstlisting}

\newpage
\chapter{Листинг обработчика измерений (pkg/benchmark/runner.go)}

В данном приложении представлен ключевой компонент, обеспечивающий воспроизводимость результатов синтетических тестов. Функция \texttt{TestRunner} выполняет разогрев, серию итераций, фильтрацию выбросов и вычисление устойчивых статистик (медиана, MAD, коэффициент вариации).

\begin{lstlisting}[language=Go, caption={Фрагмент файла pkg/benchmark/runner.go (ядро TestRunner)}]
package benchmark

import (
	"math"
	"sort"
)

type BenchmarkFunc func() (float64, string)

type BenchmarkStats struct {
	Min      float64
	Max      float64
	Avg      float64
	Median   float64
	Unit     string
	StdDev   float64
	MAD      float64
	Variance float64
}

func TestRunner(fn BenchmarkFunc, iterations int, progress chan<- float64) BenchmarkStats {
	if iterations < 10 {
		iterations = 10
	}

	// Разогрев
	_, _ = fn()
	_, _ = fn()

	results := make([]float64, iterations)
	var unit string
	for i := 0; i < iterations; i++ {
		val, u := fn()
		results[i] = val
		unit = u
		if progress != nil {
			progress <- float64(i+1) / float64(iterations)
		}
	}

	sort.Float64s(results)
	median := calculateMedian(results)
	mad := calculateMAD(results, median)
	filtered := filterOutliersIQR(results)

	var sum float64
	min := filtered[0]
	max := filtered[len(filtered)-1]
	for _, v := range filtered {
		sum += v
	}
	avg := sum / float64(len(filtered))
	filteredMedian := calculateMedian(filtered)

	var varianceSum float64
	for _, v := range filtered {
		diff := v - avg
		varianceSum += diff * diff
	}
	stdDev := math.Sqrt(varianceSum / float64(len(filtered)))
	cv := (stdDev / avg) * 100

	finalResult := 0.6*filteredMedian + 0.4*avg
	return BenchmarkStats{Min: min, Max: max, Avg: finalResult, Median: filteredMedian, Unit: unit, StdDev: stdDev, MAD: mad, Variance: cv}
}
\end{lstlisting}
